\chapter{Đánh giá Độ phức tạp Thuật toán và Kết quả Thực nghiệm}
\label{chap:chuong4}

\textit{Phụ trách biên soạn chính: Phạm Đức Cường (M5 - QA \& Kiểm thử).}

\section{Phân tích Độ phức tạp Thuật toán}
\begin{table}[htbp]
\centering
\caption{Bảng tổng hợp độ phức tạp thời gian và không gian của các giải thuật}
\label{tab:complexity}
\begin{tabular}{l c c c}
\toprule
\textbf{Thao tác} & \textbf{Trung bình} & \textbf{Xấu nhất} & \textbf{Bộ nhớ phụ trợ} \\
\midrule
Chèn sách (Insert) & $O(\log N)$ & $O(N)$ & $O(h)$ (ngăn xếp đệ quy) \\
Tìm theo ISBN (Search) & $O(\log N)$ & $O(N)$ & $O(h)$ \\
Tìm theo Tên/Tác giả & $O(N)$ & $O(N)$ & $O(h)$ \\
Xóa sách (Erase) & $O(\log N)$ & $O(N)$ & $O(h)$ \\
Duyệt In-order (LNR) & $O(N)$ & $O(N)$ & $O(h)$ \\
Giải phóng cây (Clear) & $O(N)$ & $O(N)$ & $O(h)$ \\
\bottomrule
\end{tabular}
\end{table}

\section{Danh mục Ca Kiểm thử (Test Cases)}
\section{Kết quả Chạy Thực nghiệm trên Bộ Dữ liệu Mẫu}
\section{Đánh giá Giới hạn của Cây BST và Hướng Cải tiến (Cây AVL / Red-Black Tree)}
